% =============================================================================
% ABERTURA (14:00 a 14:15)
% =============================================================================
\section[Abertura]{Abertura}

\begin{frame}{Programa do dia}
  \footnotesize
  \renewcommand{\arraystretch}{1.15}
  \begin{tabularx}{\textwidth}{@{}lX@{}}
    \toprule
    \textbf{Horário} & \textbf{Bloco} \\ \midrule
    14:00 a 14:15 & Abertura: por que avaliar? \\
    14:15 a 14:35 & O problema do contrafactual \\
    14:35 a 15:05 & Experimento controlado \hfill \destaque{Exercício 1} \\
    15:05 a 15:25 & Pareamento (PSM) \\
    15:25 a 15:55 & Diferença-em-diferenças \hfill \destaque{Exercício 2} \\
    \rowcolor{laranjaclaro}
    15:55 a 16:25 & \textbf{Intervalo} \\
    16:25 a 16:45 & Regressão descontínua \\
    16:45 a 16:55 & Outras técnicas e como escolher \\
    16:55 a 17:50 & Avaliação qualitativa e métodos mistos \hfill \destaque{Exercício 3} \\
    17:50 a 18:00 & Encerramento \\ \bottomrule
  \end{tabularx}
\end{frame}

\begin{frame}{O que vamos fazer}
  \begin{itemize}
    \item Entender por que \destaque{comparar} participantes e não participantes (ou antes e depois) \destaque{engana};
    \item Conhecer as principais técnicas quase-experimentais e o que cada uma assume;
    \item \ind{Simular em R}: sabemos a \emph{verdade} do programa e testamos se a técnica a recupera;
    \item Ver o que a avaliação \ind{qualitativa} e os \ind{métodos mistos} acrescentam aos números.
  \end{itemize}
  \vspace{0.3cm}
  \begin{block}{Como funciona}
    \small Aula expositiva + scripts em R comentados (demonstração) + 3 exercícios curtos.
  \end{block}
\end{frame}

\begin{frame}{Para começar: o programa funciona?}
  \begin{exerciciobox}{Enquete rápida (chat)}
    \small
    Um programa de \textbf{qualificação profissional} foi oferecido a trabalhadores que se inscreveram voluntariamente. Um ano depois, quem participou ganha, em média, \destaque{R\$ 587 a mais} por mês do que quem não participou.
  \end{exerciciobox}
  \vspace{0.2cm}
  O programa aumentou a renda dos participantes em R\$ 587?
  \begin{itemize}
    \item[(a)] Sim, o efeito é de R\$ 587
    \item[(b)] O efeito é menor que isso
    \item[(c)] Não dá para saber só com essa informação
  \end{itemize}
\end{frame}

\begin{frame}{A avaliação no ciclo de políticas públicas}
  \centering
  \begin{tikzpicture}[scale=0.85,every node/.style={transform shape}]
    \tikzset{est/.style={minimum width=2.6cm,minimum height=0.9cm,font=\footnotesize\bfseries}}
    \node[cxi,est] (n0) at (90:1.9cm and 1.9cm)    {Formação da\\agenda};
    \node[cxi,est] (n1) at (18:4.0cm and 1.9cm)    {Formulação};
    \node[cxi,est] (n2) at (-54:3.0cm and 1.9cm)   {Tomada de\\decisão};
    \node[cxi,est] (n3) at (-126:3.0cm and 1.9cm)  {Implementação};
    \node[cxl,est] (n4) at (162:4.0cm and 1.9cm)   {Avaliação};
    \foreach \a/\b in {0/1,1/2,2/3,3/4} { \draw[seta] (n\a) to[bend left=15] (n\b); }
    \draw[setal] (n4) to[bend left=15] (n0);
  \end{tikzpicture}
  \vspace{0.2cm}\par
  \small A avaliação \destaque{renova o ciclo}: informa se a política deve ser mantida, reformulada ou interrompida (políticas baseadas em evidências).
\end{frame}

\begin{frame}{Tipos de avaliação}
  \footnotesize
  \renewcommand{\arraystretch}{1.25}
  \begin{tabularx}{\textwidth}{@{}p{1.9cm}XX@{}}
    \toprule
    & \textbf{Pergunta} & \textbf{No programa Qualifica} \\ \midrule
    \ind{Processo} & As atividades aconteceram e os serviços chegaram ao público-alvo? & Quantos se inscreveram? Os cursos ocorreram? \\
    \ind{Resultados} & A situação do público-alvo mudou? & A renda dos participantes aumentou entre antes e depois? \\
    \rowcolor{laranjaclaro}
    \destaque{Impacto} & A mudança \textbf{é causada} pelo programa? & Quanto da variação de renda existe \emph{por causa} do treinamento? \\ \bottomrule
  \end{tabularx}
  \vspace{0.3cm}

  \small Também se classificam por \textbf{tempo} (\emph{ex-ante} e \emph{ex-post}) e por \textbf{quem avalia} (interna, externa, participativa).
\end{frame}

\begin{frame}{Cadeia de resultados (teoria do programa)}
  \centering
  \begin{tikzpicture}[scale=0.82,every node/.style={transform shape}]
    \foreach \i/\t/\ex/\sty in {0/Insumos/{Orçamento,\\instrutores}/cxi,1/Atividades/{Cursos de\\qualificação}/cxi,2/Produtos/{Trabalhadores\\treinados}/cxi,3/Resultados/{Uso das\\competências}/cxl,4/{Resultados\\finais}/{Renda e\\emprego}/cxl} {
      \node[\sty,minimum width=2.1cm,minimum height=0.9cm,font=\footnotesize\bfseries] (b\i) at (2.5*\i,0) {\t};
      \node[align=center,font=\scriptsize,text=azul] at (2.5*\i,-1.3) {\ex};
    }
    \foreach \i/\j in {0/1,1/2,2/3,3/4} { \draw[seta] (b\i) -- (b\j); }
    \draw[decorate,decoration={brace,amplitude=4pt,mirror},thick,azul] (-1,-2.1) -- (6,-2.1) node[midway,below=4pt,font=\footnotesize] {Implementação (sob controle do gestor)};
    \draw[decorate,decoration={brace,amplitude=4pt,mirror},thick,laranja] (6.2,-2.1) -- (11,-2.1) node[midway,below=4pt,font=\footnotesize] {Resultados (contexto e beneficiários)};
  \end{tikzpicture}
  \vspace{0.2cm}\par
  \small Avaliar impacto exige dizer \ind{quais resultados} o programa busca e \ind{por quais canais}.
\end{frame}
